\documentclass[11pt,a4paper]{article}
\usepackage[utf8]{inputenc}
\usepackage[T1]{fontenc}
\usepackage{lmodern}
\usepackage[margin=2.3cm]{geometry}
\usepackage{amsmath,amssymb}
\usepackage{booktabs}
\usepackage{enumitem}
\usepackage{xcolor}
\usepackage{url}
\usepackage[colorlinks=true,linkcolor=blue!50!black,citecolor=blue!50!black,urlcolor=blue!50!black]{hyperref}
\setlist{nosep}
\newcommand{\chk}{$\square$\ }

\title{Experimental protocol H4: measuring the relaxation time of a physical hyperbolic resistor--capacitor network\\ and its square control}
\author{Xavier Callens\\[2pt]
\small SocrateAI Lab, Cagnes-sur-Mer, France}
\date{Protocol version 1.0, October 2026\\ \small Companions: doi:10.5281/zenodo.23228685 (preprint v1.2), doi:10.5281/zenodo.23241463, doi:10.5281/zenodo.23244556, doi:10.5281/zenodo.23248393}

\begin{document}
\maketitle

\begin{abstract}\noindent
This protocol specifies how to build and measure two resistor--capacitor boards, a $\{7,3\}$ hyperbolic network with two layers and a square lattice disk of radius $6$, and to test three predictions fixed before any part is bought:
the ratio of their slowest relaxation times, their absolute values relative to a calibration cell, and their boundary-pair resistances. It lists the materials, the model and its numerical validation, the sessions with a gate after each, the analysis and decision rules, the failure modes with their symptoms, the requirements that
simulations place on a second, localisation phase, and the reporting rules. It is a measurement of a circuit model; it is not a test of holography and does not claim one.
\end{abstract}

\section{Purpose and predictions}
The companion preprints show by simulation that the Jacobian of the Dirichlet-to-Neumann map of a resistor network is exponentially ill-conditioned in the depth of the unknowns on flat lattices and far less so on hyperbolic tilings. Phase~1 does not measure conditioning. It checks the physical model on which all later hardware work rests:
that a real board with $1\,\%$ resistors and film capacitors relaxes with the time constants that the graph Laplacian predicts, and that the hyperbolic board is faster than the square board by the predicted factor.

\begin{center}\begin{tabular}{lll}\toprule
Id & Statement & Band \\ \midrule
T1 & $\tau(\{7,3\}\ L{=}2)\,/\,\tau(\text{square }R{=}6)$ & $0.604\pm0.03$ \\
T2 & $\tau/\tau_{\rm cal}$ for each board & $2.750\pm0.08$ and $4.551\pm0.14$ \\
S1 & six boundary-pair resistances per board (gate G2) & all within $3\,\%$ \\ \bottomrule
\end{tabular}\end{center}
The bands are frozen in the repository (preregistration 13, committed before any part was bought). You must not read the predicted values while measuring; the analysis script compares.

\paragraph{Numerical support added after the bands were fixed.}
An independent stiff integrator (the CVODE integrator of rusty-SUNDIALS) integrated the same perturbed boards that the virtual bench solves by eigendecomposition: the fitted time constants agreed to $3\times10^{-5}$, and over twelve draws of $1\,\%$ components the ratio of fitted time constants had mean $0.605$ and range $0.600$ to $0.608$.
This validates the model code, not the hardware: source resistance, probe loading, leakage and parasitics are in neither calculation.

\section{Safety and environment}
\begin{itemize}
\item Voltages are $3.3$--$5$\,V DC. The hazards are the soldering iron (burns, fumes: ventilate) and eyes (clip leads: wear glasses).
\item \textbf{Humidity} is the one environmental threat to the measurement. Surface leakage on damp perfboard is a resistance in parallel with every capacitor; with $100\,{\rm k}\Omega$ resistors a leakage path of $10\,{\rm M}\Omega$ changes $\tau$ by $1\,\%$. Store the boards in a sealed box with silica gel, measure at the same temperature for both boards, and write temperature and humidity in the records. Leakage \emph{shortens} $\tau$.
\item ESD matters for the op-amp and the converter, not for the passive parts: touch ground before handling them.
\end{itemize}

\section{Materials}
\begin{center}\begin{tabular}{p{6.8cm}p{4.4cm}p{3.6cm}}\toprule
Item & $\{7,3\}$ board & Square board \\ \midrule
Nodes & $112$ ($77$ boundary, $35$ interior) & $113$ ($69$ interior) \\
Resistors, $100\,{\rm k}\Omega$, $1\,\%$ metal film & $140$ & $200$ \\
Capacitors, $1\,\mu$F, film (PET/PP) or C0G, one per interior node & $35$ & $69$ \\
Calibration cell (one resistor, one capacitor) & $1+1$ & $1+1$ \\
Probe nodes (depths $1,2,3$) & $8,\,7,\,0$ & $10,\,20,\,31$ \\ \bottomrule
\end{tabular}\end{center}
Resistors and capacitors for both boards and both calibration cells come from the same batches. Measurement chain: a $3.3$\,V microcontroller (ESP32 or Pico) driving one rail through an op-amp follower, one unity-gain follower per probe (rail-to-rail op-amp, decoupled with $100$\,nF at each supply pin),
a 16-bit converter (ADS1115, $860$\,S/s) on I2C, USB serial to a laptop. The complete bill of materials, schematics and wiring tables are generated from the design files by the repository scripts and printed before the first session; none of them is typed by hand.
\paragraph{Why these parts.} Class-2 ceramic capacitors lose $20$--$50\,\%$ of their capacitance under bias and electrolytics leak, either of which moves $\tau$ by far more than the $3\,\%$ band. With $1\,\%$ resistors the true $\tau$ of a board is within $0.3\,\%$ of nominal; with $5\,\%$ resistors and $10\,\%$ capacitors it is within $2.7\,\%$, which uses the whole band.
Buffers are required because the converter's input would load a $100\,{\rm k}\Omega$ node by about $1\,\%$. A $10$-bit converter gave fit errors of $20$ to $80\,\%$ in simulation; $12$ bits gave at most $0.3\,\%$ and $16$ bits at most $0.1\,\%$.

\section{Model and analysis}
Each interior node $i$ has capacitance $C$ to ground; the boundary nodes are tied to one rail stepped from $0$ to $V_0=3.3$\,V at $t=0$. The interior obeys $C\,\dot V=-(L_{ii}V+L_{ib}V_b)$ with $L$ the weighted graph Laplacian. The slowest mode has $\tau=C/\lambda_{\min}(L_{ii})$, predicted as $0.275$\,s for the hyperbolic board and $0.455$\,s for the square board at the nominal values.
Analysis fits $y=1-V/V_0$ on the window $0.01<y<0.10$ by a weighted log-linear regression (weights proportional to $y$) and averages over the three probes; the window and weights were chosen on the virtual bench because the second mode decays only $2\times$ (hyperbolic) and $2.5\times$ (square) faster than the slowest, and wider windows were biased by $-2.5\,\%$ and $-6\,\%$.
The deepest probe has zero initial slope (an S-shaped onset): this is expected and is not a fault. The calibration cell gives $\tau_{\rm cal}=R_{\rm cal}C_{\rm cal}$ through the same chain, so T2 compares pure numbers.

\section{Procedure, with a gate after each session}
A gate that fails stops the sequence. You may diagnose (section~\ref{sec:fail}), repair and re-run a gate; you may not skip it.

\subsection*{Session 0: simulation and dry run (laptop, about 1\,h)}
\begin{itemize}
\item[\chk] Regenerate the design outputs, run the virtual bench, run the analysis on synthetic records (demo mode), generate the forms, print schematics, wiring tables and forms.
\item[] \textbf{Gate:} the demo prints both verdicts true. Nothing is bought before this gate.
\end{itemize}

\subsection*{Session 1: parts and sorting (multimeter, about 2\,h)}
\begin{itemize}
\item[\chk] Measure every part and write its value in the components form: resistors within $2\,\%$ of $100\,{\rm k}\Omega$, capacitors within $10\,\%$ of $1\,\mu$F. Rejected parts go to a reject pile, not back in the bag. Assign each used part to a pile (board or calibration cell).
\item[] \textbf{Gate G1:} every part used passes; the counts are $140+200+2$ resistors and $35+69+2$ capacitors.
\end{itemize}

\subsection*{Session 2: build the hyperbolic board (about 4\,h)}
\begin{itemize}
\item[\chk] Solder per the wiring table, deepest class first, one row at a time; after each row the multimeter between the row's two nodes reads the resistor (not a short, not open). One capacitor per interior node to the ground bus; boundary contacts to the rail bus; probe pads on nodes $8,7,0$; calibration cell (resistor from rail to pad, capacitor from pad to ground).
\item[\chk] Static check with the rail bus \emph{disconnected}: six boundary pairs, resistance between the two contacts.
\item[] \textbf{Gate G2:} at least five of six pairs within $3\,\%$ of the predicted resistance. A failure localises a wiring error, because the pairs share nodes.
\end{itemize}

\subsection*{Session 3: build the square board (about 5\,h)}
Same procedure; probe pads on nodes $10,20,31$; gate G2 on this board.

\subsection*{Session 4: chain and calibration (about 3\,h)}
\begin{itemize}
\item[\chk] Build the buffers; flash the logger, which writes time and three channels at about $280$\,S/s per channel for $8$\,s after the step and then holds the rail at $0$\,V.
\item[\chk] Calibration cell alone: five records per board's cell; fill the records file (file names, probe nodes, temperature, humidity).
\item[] \textbf{Gate G4:} the calibration time constant is within $5\,\%$ of $0.100$\,s and the five records agree within $1\,\%$. If not, the chain is wrong, not the boards: supply rail at the op-amp, converter gain, the cell's capacitor.
\end{itemize}

\subsection*{Session 5: measure both boards (about 2\,h, same temperature)}
\begin{itemize}
\item[\chk] Five step records per board: rail $0\to3.3$\,V, record $8$\,s, hold $0$\,V for at least $8$\,s before the next. Write temperature and humidity. Do not open the predictions file. Run the analysis.
\item[] \textbf{Gate G3:} the three probes of each record agree on $\tau$ within $3\,\%$, with at least eight samples in the fit window.
\end{itemize}

\subsection*{Session 6: record (laptop, about 1\,h)}
\begin{itemize}
\item[\chk] Write the claim for the experiment, add it to the ledger, run the experiment gate (it checks that the preregistration was committed before the data file), write the results block with numbers copied from the data file, commit the raw records, the records file, the analysis output, the components file and the session log.
\end{itemize}

\section{Decision rules}
\begin{itemize}
\item \textbf{T1 held} if the ratio lies in $0.604\pm0.03$ and gates G1--G4 passed. \textbf{T2 held} if both normalised values lie in their bands. \textbf{S1 held} if gate G2 passed on both boards.
\item A refuted prediction with all gates passed means the model of the board is wrong (ideal nodes, ideal parts, no parasitics): report it with the same prominence as a confirmation, name the suspected cause from section~\ref{sec:fail}, and leave the decision about what to change to a human.
\item A refuted prediction with a failed gate is not a result about the model.
\item Deviations from this protocol are written in the session log with the date \emph{before} the next measurement; they are never edited afterwards.
\end{itemize}

\section{What can go wrong}\label{sec:fail}
\begin{center}\small\begin{tabular}{p{4.6cm}p{5.2cm}p{4.8cm}}\toprule
Symptom & Likely cause & Check \\ \midrule
Calibration fine, board $\tau$ too \emph{long} & open resistor (a deep node isolated), or a missing capacitor's neighbour; leakage is not the cause (it shortens $\tau$) & static pairs; multimeter on the suspect row \\
Board $\tau$ too \emph{short} & leakage (humidity) or a class-2 ceramic capacitor & dry the board and remeasure; check capacitor type \\
Probes disagree by more than $3\,\%$ & a probe buffer oscillating or loading; a probe on the wrong node & scope the buffer output; continuity of the probe pad \\
S-shaped onset on the deep probe & expected & none \\
Records drift over minutes & thermal (unlikely at $1\,\%$) or incomplete discharge & wait at least $8$\,s at $0$\,V; hold temperature \\
Ratio fine, absolute $\tau$ off & calibration cell from another batch, or converter gain & remeasure the cell's resistor and capacitor \\ \bottomrule
\end{tabular}\end{center}
Further practice from the circuit literature: measure all capacitors and bin or pair them; buffers or a high-impedance probe for any scope used (a $1\,{\rm M}\Omega$ input is a $10\,\%$ load on a $100\,{\rm k}\Omega$ node); film or C0G capacitors with leakage far above $100\,{\rm k}\Omega$; shield against mains pickup; one ground point.

\section{Phase 2: what the simulations ask of the chain}
Phase~2 (not this protocol; it needs its own preregistration) measures the static Neumann-to-Dirichlet map through a $77$-channel multiplexed injection, then single-node defects and a second $\{7,3\}$ board with a switchable defect. Simulation results (not hardware) that constrain it:
\begin{itemize}
\item \textbf{Offsets are irrelevant}: any offset constant along a row or a column of the measured map is orthogonal to every defect signature, because the response matrix has zero row and column sums; the elementary linear algebra behind this is machine-checked in the repository's Lean project.
\item \textbf{Gain between the baseline and the defect map is what matters}: take both maps through the same chain in one session. Independent channel-to-channel gain differences between the two maps break localisation at about $0.3\,\%$ (square board) and $1\,\%$ ($\{7,3\}$ board); a single common gain drift is tolerated to about $1\,\%$ with the plain decoder and removed by a gain-fitted decoder.
\item \textbf{Converter resolution}: about $8$ bits for localisation on the square board and much less on the $\{7,3\}$ board; the $16$-bit converter is far above it (ideal rounding only; real converters add nonlinearity).
\item Time-resolved data do not escape the exponential growth of the conditioning with depth: in simulation, real Laplace variables lowered the condition number by about one decade and left its rate unchanged.
\item Not tested: joint worst cases with component tolerance, multi-node defects, any hardware.
\end{itemize}

\section{Data format and reporting}
The logger writes lines \texttt{t\_s,ch0,ch1,ch2}; the records file names each file with probe nodes, temperature and humidity; the analysis writes one JSON file with the three verdicts, the per-probe time constants and the gate results. Numbers go from the JSON into the results file and the ledger by copying; none is typed from memory.
Raw data are committed with the analysis output. A refuted prediction is reported as prominently as a confirmed one.

\section*{Use of AI tools}
This protocol was drafted with the assistance of Claude (Anthropic) from the repository's design files, the guide and the simulation records; the author directed the work and takes responsibility for the content.

\end{document}
